% ============================================================
%  Abstract.tex — Tóm tắt (Tiếng Việt + Tiếng Anh)
% ============================================================
\section*{TÓM TẮT}
\phantomsection
\addcontentsline{toc}{section}{\numberline{}TÓM TẮT}
\thispagestyle{empty}

Giao thức định tuyến \textbf{RPL} (IPv6 Routing Protocol for Low-Power and
Lossy Networks, RFC~6550) là chuẩn định tuyến chủ đạo trong các mạng IoT
năng lượng thấp (LLN). Tuy nhiên, thiết kế nhẹ nhàng của RPL tạo ra bề mặt
tấn công rộng mà các cơ chế bảo mật mặc định của giao thức chưa bảo vệ
được đầy đủ, đặc biệt khi đối mặt với các nút nội bộ bị xâm phạm.

Đồ án này trình bày quá trình nghiên cứu, phân tích và xây dựng \textbf{Hệ
thống Phát hiện Xâm nhập} (IDS) cho mạng RPL. Bốn dạng tấn công tiêu biểu
được tập trung nghiên cứu gồm: \textit{DIS Flooding Attack},
\textit{Blackhole Attack}, \textit{Decreased Rank Attack} và
\textit{VeRA/DODAG Version Number Attack}. Với mỗi dạng tấn công, báo cáo
phân tích cơ chế hoạt động, tác động đến hiệu năng mạng và phương pháp phát
hiện tương ứng.

Kiến trúc IDS đề xuất được tích hợp trực tiếp vào ngăn xếp Contiki-NG,
không yêu cầu sửa đổi nhân giao thức RPL. Hệ thống được kiểm thử trên bộ
mô phỏng Cooja với nhiều kịch bản khác nhau. Kết quả thực nghiệm cho thấy
IDS đề xuất đạt tỉ lệ phát hiện cao (TPR $\geq$ 90\%) với tỉ lệ cảnh báo
nhầm thấp (FPR $\leq$ 10\%), đồng thời duy trì overhead tài nguyên ở mức
chấp nhận được trên các nút cảm biến hạn chế.

\vspace{0.5cm}
\noindent\textbf{Từ khóa:} RPL, IDS, IoT, LLN, DIS Flooding, Blackhole,
Decreased Rank, VeRA, Contiki-NG, Cooja, phát hiện xâm nhập.

\vspace{1.5cm}
\begin{center}
  \rule{0.5\textwidth}{0.4pt}
\end{center}
\vspace{1cm}

\section*{ABSTRACT}
\phantomsection
\addcontentsline{toc}{section}{\numberline{}ABSTRACT}

The \textbf{RPL} protocol (IPv6 Routing Protocol for Low-Power and Lossy
Networks, RFC~6550) is the de-facto routing standard for IoT and LLN
environments. However, its lightweight design introduces a wide attack surface
that default RPL security mechanisms cannot adequately protect against,
especially when facing compromised internal nodes.

This report presents the research, analysis, and design of an \textbf{Intrusion
Detection System} (IDS) for RPL networks. Four representative attack types are
studied in depth: \textit{DIS Flooding Attack}, \textit{Blackhole Attack},
\textit{Decreased Rank Attack}, and \textit{VeRA/DODAG Version Number Attack}.
For each attack, the report analyzes the mechanism, impact on network
performance, and the corresponding detection method.

The proposed IDS architecture is integrated directly into the Contiki-NG stack
without modifying the RPL protocol core. The system is evaluated on the Cooja
Simulator under multiple scenarios. Experimental results show that the proposed
IDS achieves high detection rates (TPR $\geq$ 90\%) with low false positive
rates (FPR $\leq$ 10\%), while maintaining acceptable resource overhead on
constrained sensor nodes.

\vspace{0.5cm}
\noindent\textbf{Keywords:} RPL, IDS, IoT, LLN, DIS Flooding, Blackhole,
Decreased Rank, VeRA, Contiki-NG, Cooja, intrusion detection.

\cleardoublepage
